A directed acyclic graph consists of $n$ nodes and $m$ edges. The nodes are numbered $1,2,\dots,n$.
Calculate for each node the number of nodes you can reach from that node (including the node itself).

\vspace{0.3cm}\noindent\textbf{Input}

The first input line has two integers $n$ and $m$: the number of nodes and edges.
Then there are $m$ lines describing the edges. Each line has two distinct integers $a$ and $b$: there is an edge from node $a$ to node $b$.

\vspace{0.3cm}\noindent\textbf{Output}

Print $n$ integers: for each node the number of reachable nodes.

\vspace{0.3cm}\noindent\textbf{Constraints}

\textbullet\ $1 \le n \le 5 \cdot 10^4$
\textbullet\ $1 \le m \le 10^5$